\chapter{总结与展望}\label{chap:conclusion}

\section{总结}

\TongjiThesis{} 将《同济大学毕业设计（论文）撰写规范》转化为一套可复用的 \LaTeX{} 模板，本文档即基于该模板排版。模板的核心目标是让格式不再成为写作的负担——用户只需专注于内容，格式细节由模板处理。

在排版层面，模板复刻了官方规范的各项要求：页面版式与字体字号遵循学校规定，标题层级根据理工科/文科自动适配编号体系，图表和公式编号按章节组织，参考文献遵循 GB/T 7714 标准。在工程层面，模板基于 \pkg{ctexbook} 构建，跨平台兼容 Windows、macOS、Linux 和 Overleaf，提供 \cs{MakeCover}、\cs{MakeInfoPage}、\cs{MakeAbstract} 等语义命令简化固定页面的生成，并通过自动化持续集成确保多编译器、多平台的编译可复现。

\section{展望}

撰写规范将随学校教学管理的变化而持续演进，模板需要具备快速响应的能力。在规范层面，需密切跟踪学校对编号体系、页面版式等方面的修订，确保每届学生使用的版本始终与最新规范对齐。在工具层面，进一步降低初次使用门槛——更清晰的错误提示、更丰富的使用文档、更便捷的在线编译入口——将帮助更多同济学子从模板中受益。

模板的长久维护离不开社区共建。欢迎通过 GitHub Issues 反馈问题、通过 Pull Requests 贡献代码，共同提升模板的覆盖面与稳定性。
